\section*{Necessary group context (Section 2)}
Richard Thompson introduced three finitely presented groups $F<T<V$ acting on the unit-interval, the unit-circle and the Cantor set, respectively. A standard reference for Thompson's groups is [12]. The group $F$ consists of all homeomorphisms of the unit interval that are piecewise affine, with slopes a power of $2$ and dyadic breakpoints. Famously, $F$ has a finite presentation with two generators:
\[
F=\langle x_0,x_1\mid[x_0x_1^{-1},x_0^{-1}x_1x_0],[x_0x_1^{-1},x_0^{-2}x_1x_0^2]\rangle .
\]
Very convenient is also the following infinite presentation:
\[
F=\langle x_0,x_1,x_2,\ldots\mid x_k^{x_i}=x_{k+1}\ (i<k)\rangle .
\]
The group $F$ is orderable (so in particular torsion-free), its derived subgroup $[F,F]$ is simple and the center of $F$ is trivial.
Let $X$ be a finite set. The free monoid $X^*$ serves as the vertex set of a regular rooted tree with an edge between $v$ and $vx$ for all $v\in X^*$ and all $x\in X$. The group $W$ of automorphisms of this tree naturally acts on the set $X$ of level-1 vertices, and permutes the subtrees hanging from them. Exploiting the bijection $X^+=X^*\times X$, we thus have an isomorphism
\[
\varphi:W\to W\wr\mathrm{Sym}(X)=W^X\rtimes\mathrm{Sym}(X),
\]
mapping $g\in W$ to elements $f\in W^X$ and $\pi\in\mathrm{Sym}(X)$ as follows: $\pi$ is the restriction of $g$ to $X\subseteq X^*$, and $f$ is uniquely defined by $(xv)^g=x^\pi v^{f(x)}$. We always write $g@x$ for $f(x)$ and call it the state (or coordinate) of $g$ at $x$. If $X=\{0,\ldots,k\}$ we write $g=\langle\!\langle g@0,\ldots,g@k\rangle\!\rangle\pi$.
There exist important examples of f.g.\ self-similar weakly branched groups, notably the Grigorchuk group $G$, see [17]. It may be described as a self-similar group in the following manner: it is a group generated by $\{a,b,c,d\}$, and acts on the rooted tree $X^*$ for $X=\{0,1\}$. The action, and therefore the whole group, are defined by the restriction of $\varphi$ to $G$'s generators: $\varphi(a)=(0,1)$, $\varphi(b)=\langle\!\langle a,c\rangle\!\rangle$, $\varphi(c)=\langle\!\langle a,d\rangle\!\rangle$, and $\varphi(d)=\langle\!\langle1,b\rangle\!\rangle$, where we use the notation $(0,1)$ for the non-trivial element of $\mathrm{Sym}(X)$ (that permutes $0$ and $1$) and $\langle\!\langle w_0,w_1\rangle\!\rangle$ for a tuple in $G^{\{0,1\}}\cong G\times G$.
